All the preceding results extend directly to systems with many degrees of
freedom undergoing finite motion for which the mechanical, or classical,
problem permits complete separation of variables in the Hamilton--Jacobi
method (the so-called conditionally periodic motion; see I, §~52).  After the
variables have been separated, the problem for each degree of freedom reduces
to a one-dimensional one, and the corresponding quantization conditions are
\begin{equation}
  \oint p_i\,dq_i=2\pi\hbar(n_i+\gamma_i),
  \tag{48.6}
\end{equation}
where the integral is taken over one period of the generalized coordinate
$q_i$, while $\gamma_i$ is a number of order unity determined by the form of
the boundary conditions for that degree of freedom.\footnote{For motion in a
central field, for example,
\[
  \oint p_r\,dr=2\pi\hbar(n_r+1/2),\qquad
  \oint p_\theta\,d\theta=2\pi\hbar(l-m+1/2),\qquad
  \oint p_\varphi\,d\varphi=2\pi\hbar m
\]
(where $n_r=n-l-1$ is the radial quantum number).  The last equality follows
simply because $p_\varphi$ is the $z$ component of the angular momentum and is
equal to $\hbar m$.}
